% Ad Atticum 10.4 (ad-atticum-10-04)
% Source: https://raw.githubusercontent.com/PerseusDL/canonical-latinLit/master/data/phi0474/phi057/phi0474.phi057.perseus-lat1.xml
% Fetched by scripts/fetch_latin.py

% Dateline (Perseus): Scr. in Cumano xvii K. Mai. a. 705 (49).

\ciceroLetterOpener{CICERO ATTICO salutem}

\ciceroSection{1} multas a te accepi epistulas eodem die omnis diligenter scriptas, eam vero quae voluminis instar erat saepe legendam, sicuti facio. in qua non frustra laborem suscepisti, mihi quidem pergratum fecisti. qua re ut id, quoad licebit, id est quoad scies ubi simus, quam saepissime facias te vehementer rogo. ac deplorandi quidem, quod cotidie facimus, sit iam nobis aut finis omnino, si potest, aut moderatio quaedam, quod profecto potest. non enim iam quam dignitatem, quos honores, quem vitae statum amiserim cogito, sed quid consecutus sim, quid praestiterim, qua in laude vixerim, his denique in malis quid intersit inter me et istos quos propter omnia amisimus. hi sunt qui, nisi me civitate expulissent, obtinere se non putaverunt posse licentiam cupiditatum suarum. quorum societatis et sceleratae consensionis fides quo eruperit vides.

\ciceroSection{2} alter ardet furore et scelere nec remittit aliquid sed in dies ingravescit; modo Italia expulit, nunc alia ex parte persequi, ex alia provincia exspoliare conatur nec iam recusat sed quodam modo postulat ut, quem ad modum est, sic etiam appelletur tyrannus.

\ciceroSection{3} alter , is qui nos sibi quondam ad pedes stratos ne sublevabat quidem, qui se nihil contra huius voluntatem facere posse, elapsus e soceri manibus ac ferro bellum terra et mari comparat non iniustum ille quidem, sed cum pium tum etiam necessarium, suis tamen civibus exitiabile nisi vicerit, calamitosum etiam si vicerit.

\ciceroSection{4} horum ego summorum imperatorum non modo res gestas non antepono meis sed ne fortunam quidem ipsam; qua illi florentissima, nos duriore conflictati videmur. quis enim potest aut deserta per se patria aut oppressa beatus esse? et si, ut nos a te admonemur, recte in illis libris diximus nihil esse bonum nisi quod honestum, nihil malum nisi quod turpe sit, certe uterque istorum est miserrimus quorum utrique semper patriae salus et dignitas posterior sua dominatione et domesticis commodis fuit.

\ciceroSection{5} praeclara igitur conscientia sustentor, cum cogito me de re publica aut meruisse optime cum potuerim, aut certe numquam nisi pie cogitasse, eaque ipsa tempestate eversam esse rem publicam quam ego XIIII annis ante prospexerim. hac igitur conscientia comite proficiscar magno equidem cum dolore nec tam id propter me aut propter fratrem meum, quorum est iam acta aetas, quam propter pueros, quibus interdum videmur praestare etiam rem publicam debuisse. quorum quidem alter †non tam† quia maiore pietate est me mirabiliter excruciat, alter (o rem miseram! nihil enim mihi accidit in omni vita acerbius) indulgentia videlicet nostra depravatus eo progressus est quo non audeo dicere. et exspecto tuas litteras; scripsisti enim te scripturum esse plura cum ipsum vidisses.

\ciceroSection{6} omne meum obsequium in illum fuit cum multa severitate, neque unum eius nec parvum sed multa magna delicta compressi. Patris autem lenitas amanda potius ab illo quam tam crudeliter neglegenda. nam litteras eius ad Caesarem missas ita graviter tulimus ut te quidem celaremus sed ipsius videremur vitam insuavem reddidisse. hoc vero eius iter simulatioque pietatis qualis fuerit non audeo dicere; tantum scio post Hirtium conventum arcessitum a Caesare, cum eo de meo animo a suis rationibus alienissimo et consilio relinquendi Italiam; et haec ipsa timide. sed nulla nostra culpa est natura metuenda est. haec Curionem, haec Hortensi filium, non patrum culpa corrupit. iacet in maerore meus frater neque tam de sua vita quam de mea metuit. huic tu huic tu malo adfer consolationes, si ullas potes; maxime quidem illam velim, ea quae ad nos delata sint aut falsa esse aut minora. quae si vera sint, quid futurum sit in hac vita et fuga nescio. nam si haberemus rem publicam, consilium mihi non deesset nec ad severitatem nec ad indulgentiam. nunc haec sive iracundia sive dolore sive metu permotus gravius scripsi quam aut tuus in illum amor aut meus postulabat; si vera sunt, ignosces, sin falsa, me libente eripies mihi hunc errorem. quoquo modo vero se res habebit, nihil adsignabis nec patruo nec patri.

\ciceroSection{7} cum haec scripsissem, a Curione mihi nuntiatum est eum ad me venire. venerat enim is in Cumanum vesperi pridie, id est Idibus. si quid igitur eius modi sermo eius attulerit quod ad te scribendum sit, id his litteris adiungam.

\ciceroSection{8} praeteriit villam meam Curio iussitque mihi nuntiari mox se venturum cucurritque Puteolos ut ibi contionaretur. contionatus est, rediit, fuit ad me sane diu. o rem foedam! nosti hominem; nihil occultavit, in primis nihil esse certius quam ut omnes qui lege Pompeia condemnati essent restituerentur; itaque se in Sicilia eorum opera usurum. de Hispaniis non dubitabat quin Caesaris essent. inde ipsum cum exercitu ubicumque Pompeius esset. eius interitu finem †illi† fore. propius factum esse nihil; †ei† plane iracundia elatum voluisse Caesarem occidi Metellum tribunum pl. quod si esset factum, caedem magnam futuram fuisse. permultos hortatores esse caedis, ipsum autem non voluntate aut natura non esse crudelem, sed quod putaret popularem esse clementiam. quod si populi studium amisisset, crudelem fore. eumque perturbatum quod intellegeret se apud ipsam plebem offendisse de aerario. itaque ei cum certissimum fuisset ante quam proficisceretur contionem habere, ausum non esse vehementerque animo perturbato profectum.

\ciceroSection{9} cum autem ex eo quaererem quid videret, quem exitum, quam rem publicam, plane fatebatur nullam spem reliquam. Pompei classem timebat. quae si exisset, se de Sicilia abiturum. quid isti inquam sex tui fasces? si a senatu, cur laureati? si ab ipso, cur sex? cupivi inquit ex senatus consulto surrupto; nam aliter non poterat. at ille impendio nunc magis odit senatum. a me inquit omnia proficiscentur. cur autem sex?

\ciceroSection{10} quia xii nolui; nam licebat. tum ego quam vellem inquam petisse ab eo quod audio Philippum impetrasse! sed veritus sum, quia ille a me nihil impetrabat. libenter inquit tibi concessisset. verum puta te impetrasse; ego enim ad eum scribam, ut tu ipse voles, de ea re nos inter nos locutos. quid autem illius interest, quoniam in senatum non venis, ubi sis? quin nunc ipsum minime offendisses eius causam, si in Italia non fuisses. ad quae ego me recessum et solitudinem quaerere, maxime quod lictores haberem. laudavit consilium. quid ergo? inquam; nam mihi cursus in Graeciam per tuam provinciam est, quoniam ad mare superum milites sunt. quid mihi inquit optatius? hoc loco multa perliberaliter. ergo hoc quidem est profectum, ut non modo tuto verum etiam palam navigaremus.

\ciceroSection{11} reliqua in posterum diem distulit; ex quibus scribam ad te si quid erit epistula dignum. sunt autem quae praeterii, interregnumne esset exspectaturus, an, quo modo dixerit ille quidem, ad se deferri consulatum sed se nolle in proximum annum. et alia sunt quae exquiram. iurabat ad summam, quod nullo negotio facere solet , amicissimum mihi Caesarem esse debere. quid enim umquam? scripsit ad me Dolabella. dic quid? adfirmabat eum scripsisse, quod me cuperet ad urbem venire, illum quidem gratias agere maximas et non modo probare sed etiam gaudere. quid quaeris? acquievi. levata est enim suspicio illa domestici mali et sermonis Hirtiani. quam cupio illum dignum esse nobis et quam ipse me invito, †quae pro illo sit suspicandum!† sed opus fuit, Hirtio convento. est profecto nescio quid, sed velim quam minimo. et tamen eum nondum redisse miramur. sed haec videbimus.

\ciceroSection{12} tu Oppios Terentiae delegabis. iam enim urbis vanum periculum est. me tamen consilio iuva, pedibusne Regium an hinc statim in navem, et cetera, quoniam commoror. ego ad te statim habebo quod scribam simul et videro Curionem. de Tirone cura, quaeso, quod facis, ut sciam quid is agat.
